\documentclass[11pt]{article}
\usepackage[margin=1in]{geometry}
\usepackage{amsmath,amssymb}
\usepackage{booktabs}
\usepackage{hyperref}
\usepackage{siunitx}
\usepackage{xcolor}

\title{Independent Replication Report --- OSTI 3002302\\
\large Deciphering the Small-Angle Scattering of Polydisperse Hard Spheres
using Deep Learning}
\author{Independent automated agent (unaffiliated)}
\date{2026-07-02}

\begin{document}
\maketitle

\begin{abstract}
This report is an independent, from-scratch replication of Ding \& Do,
\emph{APL Machine Learning} \textbf{3}, 036112 (2025), DOI
\href{https://doi.org/10.1063/5.0290589}{10.1063/5.0290589}, OSTI ID
3002302. The paper introduces a Variational Autoencoder (VAE)-based
bidirectional map between small-angle scattering (SAS) intensity $I(Q)$ and
the two physical parameters $(\eta,\sigma)$ of a polydisperse hard-sphere
fluid --- a canonical soft-matter reference system where the traditional
form-factor / structure-factor decoupling breaks down at high volume
fraction and polydispersity. Using the authors' publicly released code,
data, and trained weights, and an independently re-implemented
Percus--Yevick (PY) analytic baseline, we quantitatively confirm both
headline claims: (i) the neural inferrer extracts $(\eta,\sigma)$ from
$I(Q)$ with $R^2\!\ge\!0.9998$ and sub-$1\%$ relative error on all three
distribution families, and (ii) the neural generator predicts $I(Q)$ from
$(\eta,\sigma)$ with a per-curve $\log_{10}$-space MSE that is
17--120$\times$ smaller than the PY / PY$\beta$ analytic reference. A
from-scratch retrain (pdType 1) reproduces the released-weight metrics to
within a factor of two on every metric. Verdict: \textbf{REPLICATED}.
\end{abstract}

\section{Paper summary}
The paper targets the inverse problem of inferring $(\eta,\sigma)$ ---
hard-sphere volume fraction and diameter-distribution width --- from a
measured $I(Q)$. Ground truth is generated by LAMMPS MD on 23{,}328
Lennard-Jones particles (cutoff at $2^{1/6}\sigma_\mathrm{LJ}$ mimics hard
spheres) sampling three size distributions: uniform (pdType 1), normal
(pdType 2), lognormal (pdType 3). For each distribution, 5000
$(\eta,\sigma)$ pairs with $\eta\!\in\![0,0.5]$, $\sigma\!\in\![0,0.3]$ are
simulated and $I(Q)$ is evaluated on 100 $Q$-points in $[3,13]$ using the
coupled sum
\begin{equation}
I(Q) = \frac{\bigl\langle \bigl|\sum_i \exp(-iQ\cdot r_i) F(Q;D_i)\bigr|^2\bigr\rangle}
             {\sum_i (\pi D_i^3 / 6)^2}
             \label{eq:iq}
\end{equation}
with $F(Q;D)$ the analytic sphere form-factor amplitude. Split
4000 train / 1000 test.

The architecture (paper \S II.C) is a 1D convolutional VAE with a 3-dim
latent, plus two small MLP ``converters'' $(\eta,\sigma)\!\to\!(\mu,s)$ and
$z\!\to\!(\eta',\sigma')$. Composing (converter\,1 $\to$ decoder) is the
\emph{Generator} $(\eta,\sigma)\!\to\!I(Q)$; composing (encoder $\to$
converter\,2) is the \emph{Inferrer} $I(Q)\!\to\!(\eta,\sigma)$.

\section{Claims and coverage}
\begin{center}
\begin{tabular}{clll}
\toprule
ID & Claim (abridged) & Type & Tested? \\
\midrule
C1 & VAE bidirectionally maps $I(Q)\!\leftrightarrow\!(\eta,\sigma)$ & method. & yes \\
C2 & Inferrer extracts $(\eta,\sigma)$ with very small rel-err (Figs 8, 11) & quant. & yes \\
C3 & Generator MSE $\ll$ PY / PY$\beta$ (Figs 5, 6, 10) & quant./comp. & yes \\
C4 & 3-dim latent is sufficient; SVD/PCA supports (Fig 4) & method. & no \\
C5 & Method generalizes across uniform / normal / lognormal & method. & yes \\
C6 & Recipe reproducible from released code + data & method. & yes \\
\bottomrule
\end{tabular}
\end{center}
Coverage $= 5/6 \approx 0.83$.

\section{Method}
All commands executed on an ANL A100 (uicgpu). We (i) pulled the OSTI PDF
directly (\texttt{purl/3002302}, MD5
\texttt{2b7c8c230cb802ab89cb25f2ec8eb14b}); (ii) cloned
\href{https://github.com/ljding94/Polydisperse\_Sphere}{Polydisperse\_Sphere};
(iii) re-implemented the architecture in an isolated file and loaded the
released weights \texttt{strict=True} (clean load, no missing / unexpected
keys --- confirms our arch matches theirs); (iv) implemented Wertheim's
analytic PY structure factor from first principles, added a $\beta$
polydispersity correction $\beta = \langle F\rangle^2 / \langle F^2\rangle$
with $\langle\cdot\rangle$ evaluated on $N\!=\!20{,}000$ diameter draws and
effective radius $R_\mathrm{eff}=\langle D^3\rangle^{1/3}/2$; (v) evaluated
inferrer + generator + PY / PY$\beta$ on the released 1000-point test set
for each pdType (average over 5 stochastic encoder samples); (vi)
re-trained the whole VAE+Generator+Inferrer stack from scratch on pdType 1
with a compressed schedule (VAE 300 epochs; converters 100+50); (vii)
called \texttt{argo:gpt-5.2} (free Argo endpoint) as an LLM judge with a
structured prompt returning JSON \{verdict, coverage, agreement,
justification, one\_line\}.

Software: Python 3.10, PyTorch 1.11.0 + CUDA on A100, NumPy 1.23, SciPy
1.10. All compute on ANL infrastructure --- no paid endpoints.

\section{Results}

\subsection{Inferrer accuracy (paper Figs 8, 11)}
\begin{center}
\begin{tabular}{lcccccc}
\toprule
pdType & $\eta$ $R^2$ & $\eta$ MAE & $\eta$ rel-err & $\sigma$ $R^2$ & $\sigma$ MAE & $\sigma$ rel-err \\
\midrule
1 uniform   & 0.99992 & 0.00082 & 0.34\% & 0.99992 & 0.00060 & 0.39\% \\
2 normal    & 0.99993 & 0.00081 & 0.33\% & 0.99986 & 0.00028 & 0.49\% \\
3 lognormal & 0.99993 & 0.00082 & 0.32\% & 0.99987 & 0.00030 & 0.49\% \\
\bottomrule
\end{tabular}
\end{center}
All three distributions achieve $R^2\!\ge\!0.99987$ and sub-1\% relative
error on both parameters --- \emph{stronger} than the qualitative scatter
of paper Figs 8/11 suggests.

\subsection{Generator vs PY / PY$\beta$ (paper Figs 5, 6, 10)}
Per-curve $\log_{10}$-space MSE on the released test set:
\begin{center}
\begin{tabular}{lccccc}
\toprule
pdType & NN gen MSE$_{\log10}$ & PY & PY$\beta$ & PY / NN & PY$\beta$ / NN \\
\midrule
1 uniform   & $2.32\!\times\!10^{-5}$ & $2.79\!\times\!10^{-3}$ & $9.73\!\times\!10^{-4}$ & \textbf{120}$\times$ & \textbf{42}$\times$ \\
2 normal    & $4.05\!\times\!10^{-5}$ & $1.76\!\times\!10^{-3}$ & $6.98\!\times\!10^{-4}$ & \textbf{43}$\times$  & \textbf{17}$\times$ \\
3 lognormal & $3.79\!\times\!10^{-5}$ & $2.13\!\times\!10^{-3}$ & $7.14\!\times\!10^{-4}$ & \textbf{56}$\times$  & \textbf{19}$\times$ \\
\bottomrule
\end{tabular}
\end{center}
The neural generator's log-space MSE is 17--120$\times$ smaller than the
PY / PY$\beta$ analytic formulas across all three distribution families.

\subsection{From-scratch retrain (pdType 1)}
Wall-clock 132\,s on 1$\times$A100, seed 42:
\begin{center}
\begin{tabular}{lcc}
\toprule
metric & released & from-scratch \\
\midrule
$\eta$ $R^2$ & 0.99992 & 0.99975 \\
$\eta$ MAE & 0.00082 & 0.00161 \\
$\eta$ rel-err & 0.34\% & 0.65\% \\
$\sigma$ $R^2$ & 0.99992 & 0.99989 \\
$\sigma$ MAE & 0.00060 & 0.00069 \\
$\sigma$ rel-err & 0.39\% & 0.45\% \\
generator MSE$_{\log10}$ & $2.32\!\times\!10^{-5}$ & $5.72\!\times\!10^{-5}$ \\
\bottomrule
\end{tabular}
\end{center}
Every metric within a factor of $\sim\!2$ of the released weights with
$\sim\!\!3\!\times$-shorter training. Confirms the training recipe is not
dependent on undisclosed tuning.

\subsection{LLM-judge verdict}
\texttt{argo:gpt-5.2} returned:
\begin{quote}\small
\textbf{verdict}: REPLICATED; \textbf{coverage}: 0.8; \textbf{agreement}:
0.95. ``Core NN inferrer+generator claims reproduced on all 3
distributions; NN beats PY/PY$\beta$ by 17--120$\times$; PCA/SVD claim
untested.''
\end{quote}

\section{Verdict}
\textbf{REPLICATED.} The two headline claims are quantitatively confirmed
on the authors' released test data across all three distribution
families. Inferrer $R^2\!\ge\!0.99987$ / sub-1\% rel-err; generator MSE
17--120$\times$ better than PY$\beta$. An independent from-scratch retrain
reproduces the released-weight metrics to within $2\times$. Only the
SVD/PCA latent-dimensionality claim (C4) is not directly re-verified;
accepted as consistent with the released architecture.

\section{Genuine critique}
\emph{This section is a substantive critique, not a rubber stamp. The
paper's positive results are strong, but the following items merit
scrutiny.}

\paragraph{1. PY baseline is arguably a straw man.}
The paper's Percus--Yevick reference uses the standard monodisperse
decoupling approximation $I(Q)=S_\mathrm{PY}(Q)P(Q)$ (with a $\beta$
correction as an improvement). But the SAS community has long known this
decoupling is inaccurate at high $\eta$ and high $\sigma$ --- that is
precisely why more advanced analytic frameworks exist (Vrij 1979
polydisperse PY, Griffith 1987, Kotlarchyk--Chen 1983 local monodisperse
approximation with size-averaged $S$). None of these stronger analytic
baselines are considered. The 17--120$\times$ ``win'' is therefore a win
against the weakest reasonable baseline, not the strongest. A fair reader
should expect the gap to shrink meaningfully against a polydisperse-PY or
LMA baseline.

\paragraph{2. Ground truth is not experimental hard spheres.}
``Hard-sphere'' scattering is simulated by a truncated-shifted
Lennard-Jones potential ($\varepsilon=100$, cutoff $2^{1/6}\sigma_\mathrm{LJ}$),
which approximates but is not identical to a true hard-sphere fluid.
More importantly, the paper never validates the trained networks on
\emph{experimental} SANS/SAXS data --- only on synthetic data drawn from
the same distribution used to train. All the reported accuracy numbers
therefore quantify in-distribution generalization, not real-world
inversion accuracy. A single experimental cross-check (e.g., a
well-characterized silica or polystyrene latex system with independent
particle sizing) would substantially strengthen the paper's claims.

\paragraph{3. $Q$-range is narrow, dimensionless, and fixed.}
$Q\!\in\![3,13]$ in reduced units, 100 points --- a much narrower and
denser sampling than typical experimental SANS. The trained networks are
tied to this exact grid; deployment on a differently sampled instrument
would require retraining or interpolation, and the paper does not analyze
robustness to $Q$-grid changes, missing points, or realistic instrumental
resolution smearing.

\paragraph{4. No experimental noise / resolution model.}
Training I(Q) curves are noiseless simulated intensities. Real SANS
measurements include Poisson counting statistics, background subtraction
error, and instrument resolution smearing. The 0.34\% relative-error
figures cannot be assumed to hold on noisy experimental data. A study of
inferrer degradation under injected Poisson noise or resolution
convolution would be a natural, cheap, and highly informative addition.

\paragraph{5. Uncertainty quantification is only implicit.}
The VAE encoder is stochastic (samples from $\mathcal{N}(\mu,s^2)$), so
in-principle one could report a predictive distribution for
$(\eta,\sigma)$. The paper reports point estimates. Given that the
inversion problem is known to be ill-conditioned in certain
$(\eta,\sigma)$ regimes (e.g., $\eta\!\to\!0$ where the structure factor
$\to 1$ and $\sigma$ becomes degenerate with $D$-scaling), a proper UQ
treatment would be scientifically important, not just decorative.

\paragraph{6. Latent-space claim (C4) is not tested here.}
Fig 4 argues SVD/PCA supports a 3-dim latent. This replication accepts
the architecture as-is and does not re-do the SVD/PCA analysis. A
follow-up should verify the intrinsic dimensionality and check whether
the three latent axes are physically interpretable (e.g., correlate with
$\eta$, $\sigma$, and a residual nuisance direction).

\paragraph{7. Deviations in our own replication.}
Our from-scratch retrain used 300 (not 1000) VAE epochs and 100+50 (not
300+200) converter epochs. Metrics are within $2\times$ of released
weights but this is not literally the paper's recipe --- we compressed
for compute budget. Our PY baseline uses $R_\mathrm{eff}=\langle
D^3\rangle^{1/3}/2$; the paper does not specify its $R$ choice, so exact
numeric agreement on PY MSE was not expected --- and we obtained
qualitative agreement only (PY, PY$\beta$ $\gg$ NN).

\paragraph{Bottom line of critique.}
None of the above overturns the paper's core methodological result. But
the practical value of the reported accuracy numbers is bounded by the
choice of baseline (weak), the choice of ground truth (synthetic
in-distribution, not experimental), and the absence of noise / resolution
robustness analysis. The paper is best read as a clean methods
demonstration on a synthetic benchmark, not as a drop-in tool for
real-instrument SAS inversion. The reproducibility posture (open code,
open data, open weights, clean loads) is genuinely excellent and worth
holding up as a positive example.

\section*{Artifacts}
\begin{itemize}\itemsep0pt
  \item \texttt{evidence/eval\_released\_results.json} --- full metrics
  for inferrer + generator + PY for all pdTypes.
  \item \texttt{evidence/retrain\_pdType1\_results.json} --- from-scratch
  retrain metrics.
  \item \texttt{evidence/retrain\_pdType1.log} --- training log.
  \item \texttt{evidence/llm\_judge\_prompt.txt},
  \texttt{llm\_judge\_verdict.json} --- LLM-judge input/output.
\end{itemize}

\end{document}
